% start, the call's inputs; during the call, values read from storage). Illustrative contract; the analyzer works on bytecode,
% the Solidity source is shown for the reader. (a) the code, with a fee sent to the account stored in feeTo; (b) the access
% plan: each read's location expression with its condition and its binding time; the location of balances[feeTo] depends on
% the value of feeTo, read during the call, and is not in the plan (grey italic: left to execution); (c) the pool's call in
% the swap of the walkthrough, transfer(alice, 5): the plan yields three locations at the call start. Blue bold marks the call
% inputs that the locations depend on, in the code, in the plan and in the values bound at the call start; grey italic marks
% a value read from storage.
% binding times 6.3 pt, and the expressions mixed Helvetica letters with Times digits and Computer Modern symbols); the
% expressions are written without subscripts (k1, arg0) and with || for concatenation, so every character is in one face of
% at least 8 pt; the binding times without "at"; code stays in the text's typewriter face at 8 pt; the arrows of (c) are
% drawn (the math arrow was Computer Modern).
\begin{tikzpicture}[x=1cm, y=1cm, font=\eapfont,
  pane/.style={draw=black!55, line width=0.5pt, rounded corners=2pt, inner sep=2.5pt, anchor=north west, text width=8.1cm, align=left,
    font=\fontsize{9}{9.8}\selectfont\sffamily, execute at begin node={\renewcommand\arraystretch{1.0}}},
  ttl/.style={font=\eapfont\bfseries, anchor=south west, inner sep=1pt},
]
\newcommand{\cc}[1]{{\ttfamily\footnotesize #1}}
\newcommand{\ct}[1]{{\ttfamily\mdseries\footnotesize #1}}
\newcommand{\hc}[1]{{\ttfamily\footnotesize\bfseries\color{cPlan!85!black}#1}}
\newcommand{\gc}[1]{{\ttfamily\footnotesize\slshape\color{black!55}#1}}
% plan expressions: Helvetica throughout; \hl a call input (blue bold), \cat concatenation
\newcommand{\hl}[1]{{\color{cPlan!85!black}\bfseries #1}}
\newcommand{\cat}{\,\textbar\kern0.6pt\textbar\,}
\newcommand{\bt}[1]{{\color{black!70}#1}}
% "yields": a drawn arrow (the math arrow is Computer Modern)
\newcommand{\yields}{\tikz[baseline=-0.55ex]\draw[-{Latex[length=3pt, width=3pt]}, line width=0.6pt] (0,0) -- (10pt,0);}
% (a) source
\node[ttl] at (0,0.05) {(a) Contract code};
\node[pane, fill=black!4] (src) at (0,0) {%
\begin{tabular}{@{}l@{\ \ \ }l@{}}\cc{mapping(address => uint) balances;} & \cc{// position 0}\\
\cc{mapping(address => address) referrer;} & \cc{// position 1}\end{tabular}\\
\cc{address constant TREASURY = 0x7a..;}\\
\cc{function transfer(address to, uint v) \{}\\
\cc{\ \ balances[}\hc{msg.sender}\cc{] -= v;\ \ balances[}\hc{to}\cc{] += v - v/50;}\\
\cc{\ \ balances[TREASURY] += v/100;}\\
\cc{\ \ balances[}\gc{referrer[}\hc{to}\gc{]}\cc{] += v/100; \}}\\
\cc{function airdrop(address[] r, uint v) \{}\\
\cc{\ \ for (i = 0; i < r.length; i++) balances[}\hc{r[i]}\cc{] += v; \}}};
% (b) plan
\node[ttl, anchor=north west] (tb) at ([yshift=-2pt]src.south west) {(b) Access plan, derived at compile time};
\node[pane, fill=bpOursFill] (plan) at (tb.south west) {%
\begin{tabular*}{\linewidth}{@{}l@{\ }l@{\extracolsep{\fill}}r@{}}
\cc{transfer}: & k1 = keccak(\hl{caller}\cat 0) & \bt{call start}\\
 & k2 = keccak(\hl{arg0}\cat 0) & \bt{call start}\\
 & k3 = keccak(TREASURY\cat 0) & \bt{compile time}\\
 & k4 = keccak(\hl{arg0}\cat 1) & \bt{call start}\\
 & {\color{black!55}keccak(\textit{referrer[to]}\cat 0)} & \bt{\textit{during the call}}\\
\cc{airdrop}: & k(i) = keccak(\hl{arg0}[i]\cat 0),\ i\,<\,len(\hl{arg0}) & \bt{call start}
\end{tabular*}};
% (c) instantiation
\node[ttl, anchor=north west] (tc) at ([yshift=-2pt]plan.south west) {(c) At call start: the pool calls \ct{transfer(alice, 5)}};
\node[pane, fill=white] (inst) at (tc.south west) {%
\begin{tabular}{@{}l@{\quad}l@{}}
k1 = keccak(\hl{pool}\cat 0) & \yields\ read \cc{balances[pool]}\\
k2 = keccak(\hl{alice}\cat 0) & \yields\ read \cc{balances[alice]}\\
k3 = keccak(TREASURY\cat 0) & \yields\ read \cc{balances[TREASURY]}\\
k4 = keccak(\hl{alice}\cat 1) & \yields\ read \cc{referrer[alice]}
\end{tabular}};
\end{tikzpicture}
